% GENERATED FILE. Edit canonical lessons, then run npm run generate:books.
% canonical-source-hash: fnv1a64:13c61b6fd238a54e
% canonical-lessons: LA-C431-habena, LA-C431-lorum, LA-C431-frenum, LA-C431-freno, LA-C431-scutica

\chapter{Rein, Leather strap, Bridle, To bridle, Whip}
\label{ch:la-rein-leather-strap-bridle-to-bridle-whip}
\hlchaptermodality{431}

\begin{chapteropening}
\textbf{By the end of this chapter:} \emph{I can talk about steering a horse: the reins, a leather strap, the bridle, bridling, and the whip.}

Complete the last lesson of chapter 431: I can talk about steering a horse: the reins, a leather strap, the bridle, bridling, and the whip.
\end{chapteropening}

\section[habēna, habēnae]{habēna, habēnae — a rein; a strap}

\label{lesson:LA-C431-habena}



\begin{quote}
Before the new one: say the Latin for to yoke, then the Latin for to unyoke.
\end{quote}

\subsection*{You'll want to know: habēna, habēnae}
\textbf{habēna, habēnae} — \textquotedblleft{}a rein; a strap\textquotedblright{}.

\begin{quote}
Raedārius habēnās tenet et equī lentē eunt.
\end{quote}

\begin{quote}
\emph{The coachman holds the reins and the horses go slowly.}
\end{quote}

First declension, feminine: \emph{habēna, habēnae f.}, mostly plural: \emph{habēnās dare}, to give a horse its head.

\begin{cousinweb}
\emph{habēna} is built on \emph{habeō} 'to have, to hold' (EtymDB; Lewis \& Short): the strap by which a horse is held.
\end{cousinweb}

\begin{grammarlens}[title={What you've built}]
One more word for telling, explaining and describing at B1.
\end{grammarlens}

\subsection*{Your turn}
\begin{itemize}
\raggedright
  \item \emph{Say it:} \emph{habēna, habēnae}
  \item \emph{Say it:} \emph{habēna, habēnae}, once more
  \item \emph{Say it:} a sentence of your own with \emph{habēna, habēnae}
  \item \emph{Recall:} say the Latin for to yoke, then the Latin for to unyoke, then say \emph{habēna, habēnae} again
\end{itemize}

\begin{tcolorbox}[breakable,colback=teal!4,colframe=teal!35!black,title={Before you move on}]
What does habēna, habēnae mean? (\textquotedblleft{}A rein; a strap\textquotedblright{}.) Say it once more.
\end{tcolorbox}
\section[lōrum, lōrī]{lōrum, lōrī — a leather strap, a thong; (plural) reins}

\label{lesson:LA-C431-lorum}



\begin{quote}
Before the new one: say the Latin for to unyoke, then the Latin for a rein.
\end{quote}

\subsection*{You'll want to know: lōrum, lōrī}
\textbf{lōrum, lōrī} — \textquotedblleft{}a leather strap, a thong; (plural) reins\textquotedblright{}.

\begin{quote}
Servus sarcinās lōrīs ad mūlum ligat.
\end{quote}

\begin{quote}
\emph{The slave ties the bags to the mule with straps.}
\end{quote}

Second declension, neuter: \emph{lōrum, lōrī n.} A strap of leather for tying, for reins, or for whipping.

\begin{cousinweb}
\emph{lōrum} is inherited. Lewis \& Short compare Greek \emph{eulēra} 'reins' and suggest an older form \emph{vlōrum}; EtymDB gives no parent, so the link is hedged.
\end{cousinweb}

\begin{grammarlens}[title={What you've built}]
One more word for telling, explaining and describing at B1.
\end{grammarlens}

\subsection*{Your turn}
\begin{itemize}
\raggedright
  \item \emph{Say it:} \emph{lōrum, lōrī}
  \item \emph{Say it:} \emph{lōrum, lōrī}, once more
  \item \emph{Say it:} a sentence of your own with \emph{lōrum, lōrī}
  \item \emph{Recall:} say the Latin for to unyoke, then the Latin for a rein, then say \emph{lōrum, lōrī} again
\end{itemize}

\begin{tcolorbox}[breakable,colback=teal!4,colframe=teal!35!black,title={Before you move on}]
What does lōrum, lōrī mean? (\textquotedblleft{}A leather strap, a thong; (plural) reins\textquotedblright{}.) Say it once more.
\end{tcolorbox}
\section[frēnum, frēnī]{frēnum, frēnī — a bridle, a bit}

\label{lesson:LA-C431-frenum}



\begin{quote}
Before the new one: say the Latin for a rein, then the Latin for a leather strap.
\end{quote}

\subsection*{You'll want to know: frēnum, frēnī}
\textbf{frēnum, frēnī} — \textquotedblleft{}a bridle, a bit\textquotedblright{}.

\begin{quote}
Equus ferōx frēnum nōn amat.
\end{quote}

\begin{quote}
\emph{A wild horse does not like the bridle.}
\end{quote}

Second declension, neuter: \emph{frēnum, frēnī n.}; the plural is often \emph{frēnī m.} Also spelled \emph{fraenum}. English \emph{refrain} 'to hold back' comes from \emph{refrēnāre} 'to bridle back'.

\begin{cousinweb}
\emph{frēnum} is inherited and is usually traced to PIE \emph{d\textsuperscript{h}er-} 'to hold, to support' (EtymDB; Lewis \& Short). English \emph{refrain} 'to hold back' comes from Latin \emph{refrēnāre}, built on it.
\end{cousinweb}

\begin{grammarlens}[title={What you've built}]
One more word for telling, explaining and describing at B1.
\end{grammarlens}

\subsection*{Your turn}
\begin{itemize}
\raggedright
  \item \emph{Say it:} \emph{frēnum, frēnī}
  \item \emph{Say it:} \emph{frēnum, frēnī}, once more
  \item \emph{Say it:} a sentence of your own with \emph{frēnum, frēnī}
  \item \emph{Recall:} say the Latin for a rein, then the Latin for a leather strap, then say \emph{frēnum, frēnī} again
\end{itemize}

\begin{tcolorbox}[breakable,colback=teal!4,colframe=teal!35!black,title={Before you move on}]
What does frēnum, frēnī mean? (\textquotedblleft{}A bridle, a bit\textquotedblright{}.) Say it once more.
\end{tcolorbox}
\section[frēnō, frēnāre]{frēnō, frēnāre — to bridle; to hold back}

\label{lesson:LA-C431-freno}



\begin{quote}
Before the new one: say the Latin for a leather strap, then the Latin for a bridle.
\end{quote}

\subsection*{You'll want to know: frēnō, frēnāre}
\textbf{frēnō, frēnāre} — \textquotedblleft{}to bridle; to hold back\textquotedblright{}.

\begin{quote}
Puer equum frēnat, deinde ascendit.
\end{quote}

\begin{quote}
\emph{The boy bridles the horse, then gets on.}
\end{quote}

\emph{frēnō, frēnāre, frēnāvī, frēnātum}. Also figurative: \emph{īram frēnāre}, to hold back anger.

\begin{cousinweb}
\emph{frēnō} is built on \emph{frēnum} 'bridle' (Lewis \& Short), a regular first-conjugation verb made from the noun.
\end{cousinweb}

\begin{grammarlens}[title={What you've built}]
One more word for telling, explaining and describing at B1.
\end{grammarlens}

\subsection*{Your turn}
\begin{itemize}
\raggedright
  \item \emph{Say it:} \emph{frēnō, frēnāre}
  \item \emph{Say it:} \emph{frēnō, frēnāre}, once more
  \item \emph{Say it:} a sentence of your own with \emph{frēnō, frēnāre}
  \item \emph{Recall:} say the Latin for a leather strap, then the Latin for a bridle, then say \emph{frēnō, frēnāre} again
\end{itemize}

\begin{tcolorbox}[breakable,colback=teal!4,colframe=teal!35!black,title={Before you move on}]
What does frēnō, frēnāre mean? (\textquotedblleft{}To bridle; to hold back\textquotedblright{}.) Say it once more.
\end{tcolorbox}
\section[scutica, scuticae]{scutica, scuticae — a whip, a lash}

\label{lesson:LA-C431-scutica}



\begin{quote}
Before the new one: say the Latin for a bridle, then the Latin for to bridle.
\end{quote}

\subsection*{You'll want to know: scutica, scuticae}
\textbf{scutica, scuticae} — \textquotedblleft{}a whip, a lash\textquotedblright{}.

\begin{quote}
Raedārius scuticam habet, sed equōs nōn tangit.
\end{quote}

\begin{quote}
\emph{The coachman has a whip, but he does not touch the horses.}
\end{quote}

First declension, feminine: \emph{scutica, scuticae f.} A strap whip of leather.

\begin{cousinweb}
\emph{scutica} was borrowed from Greek; EtymDB traces it to \emph{skytos} 'leather, hide'. Lewis \& Short gloss it as a lash or whip but give no origin.
\end{cousinweb}

\begin{grammarlens}[title={What you've built}]
One more word for telling, explaining and describing at B1.
\end{grammarlens}

\subsection*{Your turn}
\begin{itemize}
\raggedright
  \item \emph{Say it:} \emph{scutica, scuticae}
  \item \emph{Say it:} \emph{scutica, scuticae}, once more
  \item \emph{Say it:} a sentence of your own with \emph{scutica, scuticae}
  \item \emph{Recall:} say the Latin for a bridle, then the Latin for to bridle, then say \emph{scutica, scuticae} again
\end{itemize}

\begin{tcolorbox}[breakable,colback=teal!4,colframe=teal!35!black,title={Before you move on}]
What does scutica, scuticae mean? (\textquotedblleft{}A whip, a lash\textquotedblright{}.) Say it once more.
\end{tcolorbox}
